\chapter{Software Testing}
\label{ch:testing}

Testing the benchmark system means testing two things: that the pipeline
produces the corpus it claims to, and that the validation and mutation logic
behaves as specified. Because the system is a research pipeline rather than an
interactive application, much of its testing takes the form of the audit that
measures label quality and the semantics-preservation checks that guard the
mutators. This chapter describes the testing strategy, component testing,
pipeline testing, the audit as a validation procedure, and the
semantics-preservation regime for the mutators, with representative test cases.

%----------------------------------------------------------

\section{Testing Strategy}

The system is tested at three levels. Component testing checks each module
against known inputs. Pipeline testing runs the stages end to end and checks the
invariants that connect them. Validation testing measures the quality of the
output corpus through the independent audit. The three levels answer different
questions: component testing asks whether a module does what its specification
says, pipeline testing asks whether the modules compose without breaking an
invariant, and validation testing asks whether the labels the pipeline emits are
correct. A defect at any level is caught before the corpus is released, because
the audit re-measures label quality on fresh samples each round.

\section{Component Testing}

Each pipeline module is tested in isolation against known inputs.

\subsection{Sink-Presence Filter}

The sink-presence filter is tested for recall, for documentation suppression,
and for the proximity gates. Recall is checked by confirming that a sample
containing a known sink for its CWE is admitted. Documentation suppression is
checked by confirming that a sink token appearing only in a comment or docstring
does not cause admission, since the filter strips comments and docstrings before
matching. The proximity gate for CWE-22 is checked by confirming that a
filesystem sink without a nearby HTTP-request reference is rejected, which
removes test fixtures with hardcoded paths. The filter is also tested for the
test-file exclusion, which drops samples whose path marks them as test code for
the CWEs where test fixtures are a known noise source.

\subsection{Diff Filter}

The diff filter is tested against paired vulnerable and fixed versions. A
positive whose sink line differs between the two versions must be kept, and a
positive whose sink line is unchanged, even when other lines in the file differ,
must be rejected. Comment stripping and whitespace normalization are tested by
confirming that a change limited to a comment or to whitespace does not count as
a sink-line change, so a reformat-only commit does not pass a sample. The filter
is also tested for the undetermined case: when no fixed version is available, the
procedure returns a null verdict rather than a pass, which is what keeps the
\texttt{vudenc} samples in the lower-confidence tier instead of admitting them as
diff-validated.

\subsection{Negative Builder}

The negative builder is tested by confirming that a constructed safe twin
re-parses, that it preserves the sink name, and that the remediation neutralizes
the sink behavior. A twin produced from a \texttt{yaml.load} sink must contain
\texttt{yaml.safe\_load} and must still parse as valid Python. A twin produced
from a string-formatted query must use parameter binding rather than string
concatenation. The test confirms that the surface structure around the sink is
otherwise unchanged, which is what makes the twin a minimal pair against its
vulnerable original.

\subsection{Mutator Semantics Preservation}

Every mutator is tested for the property that defines it: that the mutated sample
re-parses and that its runtime behavior matches the original. Re-parsing is
checked for every mutated sample in the test set, so a transformation that
produces syntactically invalid Python is caught before evaluation. Behavior
preservation is checked per transformation against the Python evaluation rules.
The tests confirm the specific safety conditions each mutator depends on:

\begin{itemize}
  \item Dead-code injection inserts only side-effect-free statements
    (\texttt{if False: pass}, an unused assignment, a zero-iteration loop), so
    the function's return value cannot change.
  \item String splitting skips docstrings and type-hint strings, so
    \texttt{help()} output and runtime type introspection are unaffected, and it
    splits only the literal text between f-string substitutions.
  \item Identifier rename skips \texttt{self}, \texttt{cls}, builtins, imported
    names, names declared \texttt{global} or \texttt{nonlocal}, the function's
    own name, and any token inside a string literal, so dispatch, recursion, and
    string data are preserved.
  \item Wrapper extraction forwards arguments with \texttt{*\_a, **\_kw},
    generates an \texttt{async def} wrapper for an awaited sink, and inserts the
    helper after a docstring rather than before it.
  \item The sink-targeted and dataflow mutators are checked to leave the runtime
    value reaching the sink bitwise identical to the original argument.
\end{itemize}

A node a mutator cannot prove safe to transform is skipped, which is tested by
confirming that a rename colliding with an imported name is not applied and that
the surrounding code is left unchanged.

\section{Pipeline Testing}

Pipeline testing runs the stages end to end and checks the invariants. The split
is tested for the repo-group constraint by confirming zero cross-split
repository leakage, so no repository contributes samples to more than one split.
The corpus is tested for schema conformance by confirming that every sample
carries a CWE label, a matched sink pattern, a label provenance, a CVSS severity,
and a label-confidence tier. The mutation stage is tested by confirming that each
test sample yields the expected seven variants and that each variant re-parses.
The per-CVE sample cap is tested by confirming that no single CVE contributes
more than three near-identical files, which prevents one high-impact advisory
from dominating a class.

\section{Validation Testing through the Audit}

The audit is the system's primary validation test, and it follows the
controlled-measurement discipline of empirical software engineering
\cite{wohlin}. Three rounds of independent audit measured the false-positive
rate of the labels on disjoint CWE-stratified samples. Round~1 (77 samples) measured 52\% and surfaced six recurring failure
patterns, each of which became a fix in the next round. Round~2 held at 51\%
after the pattern tightenings, because new samples surfaced co-changed-file noise
on the classes the tightenings had not touched, which motivated the diff filter.
Round~3, on the diff-filtered data, measured 26\%. The audit is adversarial by
construction: reviewers see only the sample and its evidence, not each other's
verdicts, so a mislabeled sample is more likely to be caught than confirmed. The
falling false-positive rate across rounds is the evidence that the validation
pipeline works, and the residual 26\% is reported rather than hidden.

\section{Test Cases}

Table~\ref{tab:testcases} lists representative test cases and their outcomes.

\begin{table}[htb]
\caption{Representative test cases for the benchmark pipeline.}
\label{tab:testcases}
\centering
\begin{tabularx}{\textwidth}{p{0.5cm} X X c}
\toprule
\textbf{\#} & \textbf{Test} & \textbf{Expected result} & \textbf{Status} \\
\midrule
1 & Sink in code only & Sample admitted & Pass \\
2 & Sink in docstring only & Sample rejected & Pass \\
3 & CWE-22 sink, no HTTP source nearby & Sample rejected & Pass \\
4 & Sink in a test-file path & Sample excluded & Pass \\
5 & Sink line changed in fix & Positive kept & Pass \\
6 & Sink line unchanged in fix & Positive rejected & Pass \\
7 & Comment-only change in fix & Positive rejected & Pass \\
8 & No fixed version available & Undetermined, kept as MEDIUM tier & Pass \\
9 & Safe twin re-parses & Valid Python, sink name kept & Pass \\
10 & Mutated sample re-parses & Valid Python for all 7 variants & Pass \\
11 & Rename collides with import & Node skipped, code unchanged & Pass \\
12 & Repo-group split & Zero cross-split leakage & Pass \\
13 & Per-CVE cap & At most 3 files per CVE & Pass \\
\bottomrule
\end{tabularx}
\end{table}

\section{Summary}

The benchmark system is tested at the component level for each module, at the
pipeline level for its invariants, and at the validation level through the
three-round audit that measures label quality. The semantics-preservation checks
guard the mutators against changing program behavior, the repo-group check guards
the split against leakage, and the per-CVE cap guards a class against single-CVE
domination. The test cases in Table~\ref{tab:testcases} confirm the specified
behavior across all three levels. Chapter~\ref{ch:results} reports the
experimental results that this validated pipeline makes possible.
